% Exercises for Chapter 6 (L^p spaces, laws and densities). Statements only: solutions are kept
% privately. Throughout, (X, A, mu) is a measure space, m is Lebesgue measure on B(R), and
% (Omega, F, P) is a probability space, as in the notes.

\begin{exercise}[The essential supremum]\label{ex:ld-essup}
\easy Let $f,g:X\to\overline{\mathbb R}$ be measurable, and let
$S_f=\{M>0:\mu(\{|f|>M\})=0\}$, so that $\|f\|_\infty=\inf S_f$ (\noteref{def:ld-Linfty}). Recall from
the notes that the infimum is attained when $0<\|f\|_\infty<+\infty$, so that $|f|\le\|f\|_\infty$
a.e., and that $\|f\|_\infty=0$ if and only if $f=0$ a.e.; you may use these facts.
\begin{enumerate}[label=(\alph*)]
  \item Show that $\|f\|_\infty\le\sup_X|f|$, with equality if $\mu$ is the counting measure on a
    set $X$, or if $f$ is continuous on $[0,1]$ and $\mu=m$. Give a function on
    $(\mathbb R,\mathcal B(\mathbb R),m)$ that is finite everywhere with $\|f\|_\infty=+\infty$, and
    compute $\|\1_{\mathbb Q}\|_\infty$ for $m$.
  \item Show that $f=g$ a.e.\ implies $\|f\|_\infty=\|g\|_\infty$, and that
    $\|f+g\|_\infty\le\|f\|_\infty+\|g\|_\infty$ when $f+g$ is defined a.e.
  \item The quantity $\operatorname*{ess\,sup}f=\inf\{M\in\mathbb R:\mu(\{f>M\})=0\}$ is the
    essential supremum of $f$ itself. Compare it with $\|f\|_\infty$ for $f\equiv-1$ and for
    $f\ge0$, and show that $\|f\|_\infty=\max\big(\operatorname*{ess\,sup}f,\operatorname*{ess\,sup}(-f)\big)$
    when $\mu(X)>0$.
\end{enumerate}
\end{exercise}

\begin{exercise}[From $\|f\|_p$ to $\|f\|_\infty$]\label{ex:ld-lp-limit}
\medium Let $0<\mu(X)<+\infty$ and $f:X\to\overline{\mathbb R}$ measurable. Recall from the proof of
\noteref{prop:ld-inclusions} that $\|f\|_p\le\mu(X)^{\frac1p-\frac1q}\|f\|_q$ for $1\le p<q<\infty$
(\noteref{eq:ld-norm-comparison}), and $\|f\|_p\le\mu(X)^{1/p}\|f\|_\infty$; in particular, on a
probability space $p\mapsto\|f\|_p$ is non-decreasing on $[1,+\infty]$.
\begin{enumerate}[label=(\alph*)]
  \item Show that $\lim_{p\to\infty}\|f\|_p=\|f\|_\infty$, including when $\|f\|_\infty=+\infty$.
    \emph{Hint:} for $0<M<\|f\|_\infty$, $\|f\|_p\ge M\mu(\{|f|>M\})^{1/p}$.
  \item Compute $\|f\|_p$ on $([0,1],\mathcal B([0,1]),m)$ for $f(x)=2x$ and for the function of
    \noteref{exa:ld-essup}, and check (a) on them.
  \item Show that (a) fails on $((0,+\infty),\mathcal B,m)$ for $f=1$, but that it holds on any
    measure space if $f\in L^r$ for some $r<\infty$. \emph{Hint:}
    $\int|f|^p\,\dd\mu\le\|f\|_\infty^{p-r}\int|f|^r\,\dd\mu$ for $p>r$.
\end{enumerate}
\end{exercise}

\begin{exercise}[Powers of $x$ and the failure of the inclusions]\label{ex:ld-powers}
\medium This is a computational exercise. Let $a>0$ and $f_a(x)=x^{-a}$, with Lebesgue measure.
\begin{enumerate}[label=(\alph*)]
  \item Show that $f_a\in L^p((0,1))$ if and only if $ap<1$, that $f_a\in L^p((1,+\infty))$ if and
    only if $ap>1$, and that $f_a\notin L^p((0,+\infty))$ for every $p\ge1$. Which $f_a$ are in
    $L^\infty((0,1))$, and which in $L^\infty((1,+\infty))$?
  \item Let $1\le p<q<\infty$. Using (a), find a function in $L^q((1,+\infty))\setminus
    L^p((1,+\infty))$, and functions on $(0,+\infty)$ showing that neither of $L^p((0,+\infty))$ and
    $L^q((0,+\infty))$ contains the other. Which hypothesis of \noteref{prop:ld-inclusions} fails?
  \item On $(\mathbb N^*,\mathcal P(\mathbb N^*),\text{counting})$, show that $\ell^p\subseteq\ell^q$
    for $p<q$, with $\|a\|_q\le\|a\|_p$, and that $a_n=n^{-a}$ is in $\ell^q$ but not in $\ell^p$
    when $\frac1q<a\le\frac1p$.
  \item Write code that returns $\|x^{-a}\|_p$ on an interval $(\ell,h)\subseteq(0,+\infty)$ in
    closed form ($+\infty$ allowed). Tabulate the truncated integrals $\int_\varepsilon^1x^{-ap}\,\dd x$
    and $\int_1^Rx^{-ap}\,\dd x$ for $a=\frac13$, $p\in\{2,3,4\}$, $\varepsilon=10^{-2},\dots,10^{-8}$ and
    $R=10^2,\dots,10^6$, and explain from the table which integrals converge. Compare with
    adaptive quadrature (\texttt{scipy.integrate.quad}).
  \item Plot $p\mapsto\|f_{1/3}\|_p$ on $(0,1)$ for $p\in[1,3)$ and check numerically that it is
    non-decreasing, as \noteref{eq:ld-norm-comparison} requires on this probability space.
\end{enumerate}
\end{exercise}

\begin{exercise}[Distribution functions and quantiles]\label{ex:ld-cdf}
\medium Let $X$ be a random variable on $(\Omega,\mathcal F,\mathbb P)$ with distribution function
$F_X$ (\noteref{def:ld-cdf}). You may use the properties of measures of \cref{ch:02-measures}, the
uniqueness part of the extension theorem (\noteref{thm:mf-extension}), and that the finite disjoint
unions of intervals $(a,b]$, $(-\infty,b]$, $(a,+\infty)$ and $\mathbb R$ form an algebra that
generates $\mathcal B(\mathbb R)$.
\begin{enumerate}[label=(\alph*)]
  \item Show that $F_X$ is non-decreasing and right-continuous, with $\lim_{x\to-\infty}F_X(x)=0$ and
    $\lim_{x\to+\infty}F_X(x)=1$; that $F_X(x^-)=\mathbb P(X<x)$; and that
    $\mathbb P(X=x)=F_X(x)-F_X(x^-)$. Deduce that $F_X$ is continuous if $X$ is absolutely continuous,
    and that $F_X$ has at most countably many jumps.
  \item Let $F:\mathbb R\to[0,1]$ have the four properties of (a), and for $u\in(0,1)$ let
    $F^{-1}(u)=\inf\{x\in\mathbb R:F(x)\ge u\}$. Show that $F^{-1}(u)$ is a real number and that
    \[
      F^{-1}(u)\le x\iff u\le F(x)\qquad(u\in(0,1),\ x\in\mathbb R).
    \]
  \item On $((0,1),\mathcal B((0,1)),m)$, let $Y=F^{-1}$. Show that $Y$ is a random variable whose
    law, the image measure $m_{F^{-1}}$, has distribution function $F$. Conclude that every
    function with the properties of (a) is the distribution function of some random variable.
  \item Show that two random variables, possibly on different probability spaces, have the same
    law if and only if they have the same distribution function. Compute $F^{-1}$ for the law of a
    fair die and for the exponential law $F(x)=(1-e^{-\lambda x})\1_{[0,+\infty)}(x)$.
\end{enumerate}
\end{exercise}

\begin{exercise}[How many loaves to bake]\label{ex:ld-newsvendor}
\medium \emph{(In the style of Ross, Ch.~4.)} Each morning a bakery bakes $k$ loaves, at a cost of
$c$ per loaf, and sells them at a price $s>c$; loaves left in the evening are given away. The
number $D$ of customers who want a loaf on a given day is a random variable with values in
$\mathbb N$ and distribution function $F_D$, and the profit is $G_k=s\min(D,k)-ck$.
\begin{enumerate}[label=(\alph*)]
  \item Show that $\min(D,k)=\sum_{j=0}^{k-1}\1_{\{D>j\}}$, and deduce that
    $\mathbb E[\min(D,k)]=\sum_{j=0}^{k-1}\big(1-F_D(j)\big)$.
  \item Show that $\mathbb E[G_{k+1}]-\mathbb E[G_k]=s\big(1-F_D(k)\big)-c$. Deduce that the expected
    profit is largest for $k^*=\min\{k\in\mathbb N:F_D(k)\ge1-\frac cs\}$.
  \item Let $D\sim\mathcal P(20)$ (Poisson), $s=3$ and $c=1$. Compute $k^*$, $\mathbb E[G_{k^*}]$, and
    the expected profit of the naive choice $k=\mathbb E[D]=20$, numerically with the discrete density
    of $D$. Why is $k^*$ larger than the mean demand?
  \item Suppose that the loaves left can be sold to a farm at $r$ per loaf, $0\le r<c$. Find the new
    optimal number of loaves, and compute it for $r=0.5$.
\end{enumerate}
\end{exercise}

\begin{exercise}[The empirical distribution function and the sample size]\label{ex:ld-ecdf}
\hard This is a computational exercise. Let $X_1,\dots,X_n$ be independent random variables
(\cref{ch:07-product-measures}) with the same distribution function $F$, and
$F_n(x)=\frac1n\#\{k\le n:X_k\le x\}$ their empirical distribution function. You may use that the
events $\{X_k\le x\}$, $k\le n$, are independent, so that $n\,F_n(x)$ has the binomial law
$B(n,F(x))$; Chebyshev's inequality
$\mathbb P(|Y-\mathbb E[Y]|\ge\varepsilon)\le\operatorname{var}(Y)/\varepsilon^2$, and the
Dvoretzky--Kiefer--Wolfowitz inequality in the form proved by P.~Massart (\emph{Ann.\ Probab.}\ 18,
1990, 1269--1283):
$\mathbb P\big(\sup_x|F_n(x)-F(x)|>\varepsilon\big)\le2e^{-2n\varepsilon^2}$ for every $\varepsilon>0$.
\begin{enumerate}[label=(\alph*)]
  \item Write vectorised code for $F_n$, and for $D_n=\sup_x|F_n(x)-F(x)|$ computed exactly from the
    sorted sample (allow $F$ to have atoms: you will need $F(x^-)$). Check it against a fine grid.
  \item Fix $x$. Show that $\mathbb E[F_n(x)]=F(x)$ and $\operatorname{var}(F_n(x))=F(x)(1-F(x))/n\le\frac1{4n}$.
    Using Chebyshev, find $n$ such that $\mathbb P(|F_n(x)-F(x)|\ge0.01)\le0.05$ whatever $F$ and $x$.
  \item Using the DKW inequality, find $n$ such that $\mathbb P(D_n>0.01)\le0.05$. Compare with (b):
    which guarantee is stronger, and which needs more samples? What does this say about Chebyshev's
    inequality?
  \item With the $n$ of (c) and seed $2026$, draw $500$ samples of size $n$ from the exponential law
    $\mathrm{Exp}(1)$ and report the fraction with $D_n>0.01$. Repeat with the Poisson law
    $\mathcal P(3)$. Compare with $2e^{-2n\varepsilon^2}$ and with the limit
    $2\sum_{j\ge1}(-1)^{j-1}e^{-2j^2\lambda^2}$, $\lambda=\varepsilon\sqrt n$, of
    $\mathbb P(\sqrt nD_n>\lambda)$ for continuous $F$ (Kolmogorov). Is the DKW bound sharp?
\end{enumerate}
\end{exercise}

\begin{exercise}[Densities: the steps left to the reader]\label{ex:ld-density-integral}
\medium This exercise completes the proof of \noteref{prop:ld-density-integral}. Let $f\ge0$ be
measurable on $(X,\mathcal A,\mu)$, and $\nu(A)=\int_Af\,\dd\mu$ (\noteref{prop:ld-density-measure}). No
$\sigma$-finiteness is assumed in (a)--(c).
\begin{enumerate}[label=(\alph*)]
  \item Show that $\int_Xg\,\dd\nu=\int_Xgf\,\dd\mu$ in $[0,+\infty]$ for every measurable $g\ge0$.
    Deduce that $g\in L^1(\nu)$ if and only if $gf\in L^1(\mu)$, and that then
    $\int_Xg\,\dd\nu=\int_Xgf\,\dd\mu$. Deduce \noteref{cor:ld-image-density}.
  \item (Chain rule.) Let $h\ge0$ be measurable and $\lambda(A)=\int_Ah\,\dd\nu$. Show that
    $\lambda(A)=\int_Ahf\,\dd\mu$: in the notation of Radon--Nikodym derivatives,
    $\frac{\dd\lambda}{\dd\mu}=\frac{\dd\lambda}{\dd\nu}\frac{\dd\nu}{\dd\mu}$.
  \item Show that $\mu\ll\nu$ if and only if $\mu(\{f=0\})=0$, and that then $\mu(A)=\int_A\frac1f\,\dd\nu$
    for every $A\in\mathcal A$: if $\nu\sim\mu$, then $\frac{\dd\mu}{\dd\nu}=1/\frac{\dd\nu}{\dd\mu}$.
  \item Let $\mu$ be the counting measure on $([0,1],\mathcal B([0,1]))$. Show that $m\ll\mu$, but that
    there is no measurable $f\ge0$ with $m(A)=\int_Af\,\dd\mu$ for all $A$. Which hypothesis of
    \noteref{thm:ld-radon-nikodym} fails?
\end{enumerate}
\end{exercise}

\begin{exercise}[Radon--Nikodym derivatives between discrete laws]\label{ex:ld-rn-finite}
\medium This is a computational exercise. Let $\nu=B(10,0.3)$ (binomial) and $\mu=\mathcal P(3)$
(Poisson), as measures on $(\mathbb N,\mathcal P(\mathbb N))$.
\begin{enumerate}[label=(\alph*)]
  \item Show that $\nu\ll\mu$ but $\mu\not\ll\nu$, and compute $\frac{\dd\nu}{\dd\mu}(k)$ in closed form.
    Where is it largest? Check numerically that $\int g\,\dd\nu=\int g\frac{\dd\nu}{\dd\mu}\,\dd\mu$ for
    $g(k)=k$, $k^2$, $e^{k/2}$, and compare with the exact values $3$, $11.1$ and
    $(0.7+0.3e^{1/2})^{10}$.
  \item Give two laws on $\mathbb N$ that are singular, and two that are equivalent but different.
    For $0<p,q<1$, compute $\frac{\dd B(n,p)}{\dd B(n,q)}$.
  \item (Importance sampling.) The probability $\pi=\mathbb P(K\ge25)$ for $K\sim B(100,0.1)$ is
    small. Show that $\pi=\mathbb E[w(K')\1_{\{K'\ge25\}}]$ for $K'\sim B(100,0.25)$, with
    $w=\frac{\dd B(100,0.1)}{\dd B(100,0.25)}$. With seed $2026$ and $10^4$ draws, estimate $\pi$ both
    naively (draws of $K$) and with the weight (draws of $K'$), with standard errors, and compare
    with the exact value.
  \item Why would the method of (c) fail if the roles of the two laws were exchanged in a problem
    where $\nu\not\ll\mu$?
\end{enumerate}
\end{exercise}

\begin{exercise}[The area of a tile]\label{ex:ld-tile}
\easy \emph{(In the style of Ross, Ch.~5.)} This is also a computational exercise. The side $S$ of a
square tile, in centimetres, is a random variable with density
$f_S(s)=\frac34\big(1-(s-10)^2\big)\1_{(9,11)}(s)$.
\begin{enumerate}[label=(\alph*)]
  \item Check that $f_S$ is a density and compute $F_S$. Compute $\mathbb E[A]$ for the area $A=S^2$,
    using \noteref{eq:ld-ac-expectation} (no law of $A$ needed), and $\operatorname{var}(A)$.
  \item Compute $F_A$, show that $A$ is absolutely continuous, and find its density $f_A$. You may use
    the fundamental theorem of calculus and (d) of \cref{ex:ld-cdf}. Compute $\mathbb E[A]=\int af_A(a)\,\dd a$ again, and
    $\mathbb P(A>110)$.
  \item One can show that $S=10+M$, where $M$ is the median of three independent uniform random
    variables on $(-1,1)$, has density $f_S$. With seed $2026$, draw $10^6$ such $S$, check the
    claim by the distance $\sup_s|F_n(s)-F_S(s)|$, and estimate $\mathbb E[A]$ and $\mathbb P(A>110)$
    with $95\%$ intervals.
\end{enumerate}
\end{exercise}

\begin{exercise}[A deductible and a limit]\label{ex:ld-insurance}
\hard \emph{(In the style of Ross, Ch.~5.)} A loss $L$ has the exponential law with mean $1000$,
$f_L(\ell)=\lambda e^{-\lambda\ell}\1_{(0,+\infty)}(\ell)$ with $\lambda=10^{-3}$. An insurance policy
pays nothing for losses below the deductible $d=500$, pays $L-d$ above it, but never more than
$u=2500$: the payment is $Y=\min\big((L-d)^+,u\big)$.
\begin{enumerate}[label=(\alph*)]
  \item Compute $F_Y$. Show that $\mathbb P_Y$ has atoms at $0$ and at $u$, with masses
    $1-e^{-1/2}$ and $e^{-3}$, and a density on $(0,u)$, so that $Y$ is neither discrete nor
    absolutely continuous. Write $\mathbb P_Y=\mathbb P_Y^{d}+\mathbb P_Y^{ac}$, with
    $\mathbb P_Y^{d}$ concentrated on $\{0,u\}$ and $\mathbb P_Y^{ac}\ll m$.
  \item Compute $\mathbb E[Y]$ in three ways: as $\int h(\ell)f_L(\ell)\,\dd\ell$ with
    $h(\ell)=\min((\ell-d)^+,u)$; as $\int y\,\dd\mathbb P_Y(y)$, using the decomposition of (a); and as
    $\int_0^{+\infty}\mathbb P(Y>t)\,\dd t$ (show first that $\mathbb E[Z]=\int_0^{+\infty}\mathbb P(Z>t)\,\dd t$
    for every random variable $Z$ with values in $[0,u]$). Compute also $\operatorname{var}(Y)$.
  \item Where is $F_Y$ differentiable? Compute $\int_{\mathbb R}F_Y'(y)\,\dd y$, and explain why it is not
    $1$, in the light of \noteref{thm:ld-lebesgue}.
  \item With seed $2026$, simulate $10^6$ losses, and compute $95\%$ intervals for $\mathbb E[Y]$,
    $\mathbb P(Y=0)$ and $\mathbb P(Y=u)$ from the same draws. Do all three intervals contain the exact
    values? Should they?
\end{enumerate}
\end{exercise}
